% =====================================================================
%  Supporting Information scaffold — JGR:Planets
%  Purpose: receive material moved from the main text to bring the
%  manuscript inside AGU's 25-publication-unit guideline.
%
%  CURRENT PU ACCOUNTING (2026-07-04, certified build):
%    main text ~9,100 words       = 18.2 PU
%    11 figures + 4 tables        = 15   PU
%    TOTAL                        ~ 33   PU   (guideline: 25)
%
%  RECOMMENDED MOVES (need co-author sign-off; each -1 PU unless noted):
%    1. Symbols table (back matter)            -> Table S1
%    2. Fig. amplitude-vs-depth (borestem diagnostic; caveat (v)
%       already labels it diagnostic-only)     -> Figure S1
%    3. Fig. kd_qb_posterior corner panels     -> Figure S2
%       (keep fig_posterior_compare in main)
%    4. Discussion subsec. "Physical plausibility" trimmed by ~600
%       words with details here                -> Text S1 (-1.2 PU)
%    5. Discussion subsec. "Comparison with prior estimates" keeps the
%       new lineage figure + 1 paragraph; move per-study detail here
%       -> Text S2 (-1 PU)
%  Result: ~25-26 PU.
% =====================================================================
\documentclass[11pt]{article}
\usepackage[a4paper,margin=2.2cm]{geometry}
\usepackage{amsmath,amssymb,graphicx,booktabs,siunitx}
\usepackage[round,authoryear]{natbib}
\bibliographystyle{plainnat}

% AGU Supporting Information numbers figures/tables S1, S2, ...
\renewcommand{\thefigure}{S\arabic{figure}}
\renewcommand{\thetable}{S\arabic{table}}

\title{Supporting Information for\\
``Per-site meter-scale regolith thermal conductivity at the Apollo 15
and 17 boreholes''}
\author{}
\date{}

\begin{document}
\maketitle

\noindent\textbf{Contents of this file:} Text S1--S2, Figures S1--S2,
Table S1 (to be populated on co-author approval of the moves listed in
the header comment of this file).

\section*{Introduction}
This Supporting Information contains extended methodological detail and
diagnostic material referenced from the main text. Figure~S1 is the
per-site density-scalar diagnostic under the \citet{martinez2021}
$K(T,\rho)$ form, cited from \S4.3 of the main text.

% =====================================================================
\begin{figure}[!htbp]
\centering
\includegraphics[width=0.9\linewidth]{figures/fig_alpha_sweep.pdf}
\caption{Diagnostic per-site retrieval under the \citet{martinez2021}
$K(T,\rho)$ conductivity form. Meter-scale-sensor RMSE as a function of
the per-site density scalar $\alpha$ (main-text eq.~9), for both sites,
with the retrieved $\alpha^{*}$ starred. The scalar $\alpha$ rescales the
asymptotic (meter-scale) bulk density entering the Mart\'inez model,
$\rho_d^{\text{site}}=\alpha\cdot 1800\ \mathrm{kg\,m^{-3}}$, leaving the
four Mart\'inez coefficients fixed. Dashed teal: published Mart\'inez
baseline ($\alpha=1$, $\rho_d=1800\ \mathrm{kg\,m^{-3}}$). Green
verticals: Apollo-core admissible density range
($\rho_d=1700$--$2000\ \mathrm{kg\,m^{-3}}$;
\citealp{mitchell1973,carrier1991}). Dotted gray: solid-basalt upper
bound ($3000\ \mathrm{kg\,m^{-3}}$; \citealp{wieczorek2006}). Both
RMSE($\alpha$) bowls bottom inside or at the edge of the admissible band
($\alpha^{*}_{\mathrm{A15}}=1.19$, $\rho_d^{*}\approx2139\ \mathrm{kg\,m^{-3}}$;
$\alpha^{*}_{\mathrm{A17}}=1.04$, $\rho_d^{*}\approx1877\ \mathrm{kg\,m^{-3}}$),
so the published $K(T,\rho)$ is physically admissible at both boreholes.
The flat A17 bowl also shows that no admissible density adjustment closes
the residual factor-of-$\sim$1.5 RMSE gap to the per-site Hayne fit;
sharpening the A17 gradient further requires freedom in the conductivity
mechanism (the contact/radiative split), not in the density input.}
\label{fig:si-alpha-sweep}
\end{figure}

% --- remaining placeholders (populate on co-author approval) ----------
% \section*{Text S1: Physical plausibility, extended}
% \section*{Text S2: Prior estimates, study-by-study}
% \section*{Figure S2}  \includegraphics{figures/fig_kd_qb_posterior.pdf}
% \section*{Table S1: Symbols}

\bibliography{references}

\end{document}
